\section{Finite two-colour parameter transgression and cap independence}
\label{sec:r11-cap-parameter}

\paragraph{Status and scope.}
This is a new R11 construction by \texttt{/root/r11\_cube}, submitted
for two independent exact-source reviews.  It proves a finite algebraic
transgression criterion and applies it to the two prescribed capped-CS
triples for the \emph{same} square and the \emph{same} passive block.
It does not identify the six different passive-block triples in the
raw three-distinct-entry cube.  General strong G, finite-cover frame
coherence, and the entire lower-middle OPEN-KERNEL remain open.

\subsection{The finite polynomial extension of the actual algebra}

Let $W$ be an ordinary analytic open over the original $K$.  Write
$\mathscr B(W)=\bigoplus_{p,q}B^{p,q}(W)$ for the actual algebra,
with full anticommuting differentials $d',d''$ and $D=d'd''$.
Tensor it algebraically over $\mathbb R$ with
\[
 \mathbb R[t]\otimes\Lambda(\epsilon',\epsilon''),
 \qquad \deg t=(0,0),\quad
 \deg\epsilon'=(1,0),\quad\deg\epsilon''=(0,1).
 \tag{CT1}
\]
The two exterior variables anticommute with each other and with
odd-total-degree actual forms.  Extend the full differentials by
\[
 \widetilde d't=\epsilon',\qquad
 \widetilde d''t=\epsilon'',\qquad
 \widetilde d'\epsilon'=\widetilde d'\epsilon''
 =\widetilde d''\epsilon'=\widetilde d''\epsilon''=0.
 \tag{CT2}
\]
Their restrictions to $\mathscr B(W)$ are $d',d''$.  The graded
Leibniz rules imply $\widetilde d'{}^2=\widetilde d''{}^2=0$ and
$\widetilde d'\widetilde d''=-\widetilde d''\widetilde d'$.
This is a polynomial coefficient device, not an analytic parameter
space and not a pullback or pushforward involving a real interval.

Every element has a unique expression
\[
 F=F_0(t)+\epsilon'F_1(t)+\epsilon''F_2(t)
                          +\epsilon'\epsilon''F_{12}(t),
 \tag{CT3}
\]
with the exterior variables on the left and all $F_I(t)$ finite
polynomials with actual coefficients.  Evaluation
$\operatorname{ev}_a F=F_0(a)$ is a morphism of actual differential
algebras.  Coefficient integration means the finite operation
$\int_0^1\sum_jF_jt^j\,dt=\sum_jF_j/(j+1)$ and therefore lands in
actual $B$ itself.

Actual zero relations have already been imposed in $\mathscr B(W)$.
Consequently (CT1), coefficient extraction, restriction, and integration
respect all actual presentations and refinements.  On each fixed
polynomial degree this is a finite direct sum of the inherited actual
coefficient-final topologies.  All integration and extraction maps
are finite linear maps.  Products and differentials in the statements
below use one common finite refinement of the whole supplied
presentation; their finite-stage estimates are the usual full mixed
Leibniz estimates.  No global wedge-continuity theorem, strictness,
Hausdorff quotient, completion, or whole-kernel comparison is asserted.

Use the frozen pair/triple/quadruple maps
\[
 \begin{split}
 \partial_2(A,B,M)&=(d'A+d''M,\ d''B+d'M),\\
 \partial_3(U,V,P,Q)&=(d'U+d''P,\ d''V+d'Q,\ d'P+d''Q),
 \end{split}
 \tag{CT4}
\]
with $L_2^s=B^{s-2,s}\oplus B^{s,s-2}\oplus B^{s-1,s-1}$ and
$L_3^s=B^{s-3,s}\oplus B^{s,s-3}\oplus B^{s-2,s-1}
\oplus B^{s-1,s-2}$.  Negative bidegrees are zero.  Tildes mean
that the maps are formed with the differentials in (CT2).

\begin{theorem}[A finite two-colour transgression criterion]
Let $\widetilde q=(\widetilde A,\widetilde B,\widetilde M)$ be
a finite polynomial triple of total bidegrees $L_2^s$ in (CT1).
Suppose
\[
       \widetilde\partial_2\widetilde q=p,
       \qquad p\in L_1^s(W)\subset(\mathrm{CT1}),
 \tag{CT5}
\]
so the right side has neither parameter nor exterior-variable terms.
Expand its three components as in (CT3), and suppose
\[
                    A_2(0)=M_1(0)=0.
 \tag{CT6}
\]
Put
\[
 \begin{gathered}
 J=\int_0^1(1-t)M_{12}(t)\,dt,\\
 \mathcal T(\widetilde q)=
 \left(\int_0^1 A_1\,dt,\ \int_0^1 B_2\,dt,\
       \int_0^1 M_1\,dt-d''J,\ \int_0^1 M_2\,dt\right).
 \end{gathered}
 \tag{CT7}
\]
Then $\mathcal T(\widetilde q)\in L_3^s(W)$ is an explicitly
specified actual quadruple and
\[
       \operatorname{ev}_1\widetilde q
        -\operatorname{ev}_0\widetilde q
                    =\partial_3\mathcal T(\widetilde q).
 \tag{CT8}
\]
In particular this criterion does not assume or prove
$\ker\partial_2=\operatorname{im}\partial_3$.
\end{theorem}

\begin{proof}
Put a dot for the derivative of a polynomial in $t$.  The first,
second, and mixed exterior coefficients of the first component of
(CT5) give
\[
 \begin{split}
 \dot A_0&=d'A_1+d''M_1,\\
 \dot M_0&=d'A_2+d''M_2,\\
 \partial_t(A_2-M_1)&=-d'A_{12}-d''M_{12}.
 \end{split}
 \tag{CT9}
\]
The same coefficients of its second component give
\[
 \begin{split}
 \dot M_0&=d''B_1+d'M_1,\\
 \dot B_0&=d''B_2+d'M_2,\\
 \partial_t(M_2-B_1)&=-d''B_{12}-d'M_{12}.
 \end{split}
 \tag{CT10}
\]
The signs follow by moving both exterior variables to the left.
For example,
$\widetilde d'(\epsilon''A_2)
 =-\epsilon''d'A_2+\epsilon'\epsilon''\dot A_2$, whereas
$\widetilde d''(\epsilon'M_1)
 =-\epsilon'd''M_1-\epsilon'\epsilon''\dot M_1$.

The first and second components of (CT8) now follow by integrating
the first equation of (CT9) and the second of (CT10), since
$d''d''J=0$.  For the third component define
\[
 E(t)=\dot M_0-d'M_1-d''M_2=d'(A_2-M_1).
 \tag{CT11}
\]
Equation (CT9) gives $\dot E=-DM_{12}$, and (CT6) gives $E(0)=0$.
Finite polynomial integration therefore gives
\[
 \int_0^1E(t)\,dt
    =-D\int_0^1(1-t)M_{12}(t)\,dt=-DJ.
 \tag{CT12}
\]
This is exactly the correction supplied by the term $-d''J$
in the third entry of (CT7).  Thus all three components agree.
The bidegree of $J$ is $(s-2,s-2)$, so the displayed correction
has the required bidegree.  All integrations were finite scalar
coefficient integrations, proving the actual-form assertion.
\end{proof}

\subsection{Applying the criterion to two caps of the same square}

Retain exactly the data and the primitive convention of the recovered
R10 module CS.  Thus $f=(f_1,\ldots,f_r)$, $r\ge3$, is a tuple
of sections of the same metrized line bundle on $N$; $a=v(\lambda)$
and $b=v(\mu)$ are valuations of units on all of $N$; the two
changes affect entries $1,2$; $W=N\setminus|f=0|$; and $s=r-1$.
The original curvature is $R$, the four norm potentials are
$\Phi_0,\Phi_A,\Phi_B,\Phi_{AB}$, and the original projective
curvatures are $C_i=R-D\Phi_i$.

For two positive real constants $\delta_0,\delta_1$, let
\[
 \begin{split}
 \psi&=-\log\|(f_3,\ldots,f_r)\|,\\
 \Psi_i&=\min\{\psi,\Phi_0+|a|+|b|+\delta_i\},\qquad i=0,1.
 \end{split}
 \tag{CT13}
\]
At passive-zero points the first term is interpreted as $+\infty$.
The caps are actual on all of $W$ by the ordinary-open argument
in CS23--CS25.  Let $q_s[\Psi_i]$ denote \emph{exactly} CS36 with
that cap, including its entire CS30 tetrahedron expression, MT15
metric correction, scalar correction, and CPT curvature correction.

\begin{theorem}[Actual independence of the cap constant]
For the fixed whole square and fixed passive block in (CT13), there
is the following explicit actual quadruple on all of the original $W$:
\[
       H_{01}=\mathcal T(\widetilde q),\qquad
       q_s[\Psi_1]-q_s[\Psi_0]=\partial_3H_{01}.
 \tag{CT14}
\]
Here $\widetilde q$ is obtained by the following finite recipe:
\[
 \rho(t)=3t^2-2t^3,\quad
 H=\Psi_1-\Psi_0,\quad
 \widetilde\Psi=\Psi_0+\rho(t)H,\quad
 \widetilde R_*=R-\widetilde D\widetilde\Psi,\quad
 \widetilde p_i=\Phi_i-\widetilde\Psi.
 \tag{CT15}
\]
Insert these data into every term of CS36 and use the full
operators (CT2), including all $\epsilon'\epsilon''$ coefficients.
Then extract the coefficients in (CT7).  No extending error is
needed for this same-passive-block comparison.
\end{theorem}

\begin{proof}
The polynomials in (CT15) and all operations in CS36 are finite;
in particular neither an infinite series nor a parameter-dependent
choice of actual source layer occurs.  The curvature identity is
literal in (CT1):
\[
             \widetilde R_*-\widetilde D\widetilde p_i=C_i.
 \tag{CT16}
\]
All original $a,b,R,C_i$ are constant in the formal parameter.
We prove the stronger identity required by (CT5):
\[
        \widetilde\partial_2\widetilde q
                      =P_s-e_{a,b}R^{s-1},\qquad
        e_{a,b}=(a\,d''b,b\,d'a).
 \tag{CT17}
\]

First, CS10, MT8--MT21, CS32--CS35 are identities of finite
graded differential expressions.  They remain valid in (CT1).
The support step used in CS10 is unchanged: on a carrier of
$D|c|$, the original edge difference $\Phi_j-\Phi_i$ is zero
with all tangential first jets, while $|c|$ and its tangential
first jets vanish.  The normalized difference is still precisely
that same original difference by (CT15).  The exterior coefficients
therefore satisfy the identical vanishing; no mixed formal derivative
has been discarded.  Full nonclosed coefficient remainders in MT12
are retained and telescope exactly as in MT20--MT21.

It remains to verify the support identity CS29 for the formal
normalization.  The two ordinary opens
\[
 \begin{split}
 W_0&=\{\psi<\Phi_0+|a|+|b|+\min(\delta_0,\delta_1)\},\\
 W_2&=\{\psi>\Phi_0+|a|+|b|\}
 \end{split}
 \tag{CT18}
\]
cover $W$, including passive-zero points in $W_2$.  On $W_0$,
$\Psi_0=\Psi_1=\psi$, so $H=0$ as an actual scalar on this
ordinary open.  Consequently all its actual and formal derivatives
vanish there.  The formal normalization is the unchanged passive
norm normalization, and CS21 proves both sides of CS29 are zero.
This uses the original-$K$ passive projective map to
$\mathbf P^{r-3}$; no new formal-parameter dimension argument
is made.

On $W_2$, all four norms are determined by the two affected entries.
CS17--CS20 work for any common reference normalization.  To verify
this in (CT1), observe that on the old corner the relation
$\widetilde p=\widetilde q$ is the original relation $\Phi_i=\Phi_j$;
on the changed-entry new corner,
$\widetilde p+c=\widetilde q$ is the original relation
$\Phi_i+c=\Phi_j$.  Their actual tangential derivatives agree,
and their formal parameter derivatives agree because the same
$\widetilde\Psi$ is subtracted from both sides.  These facts kill
the specific symbol coefficients in CS19 before arbitrary remaining
closed curvature factors are wedged.  The rank-one wall check
$c=0$ is unchanged.  Finally the degree-one square sum is invariant
under an arbitrary common change $p_i\mapsto p_i+h$: each edge
residual changes by $\mathfrak a(h,\Phi_j-\Phi_i-c_e;1)$,
whose sum is zero by endpoint and unit-weight telescoping.
Thus CS20, and then CS18, hold coefficientwise in (CT1).

Actual sheaf locality applied to each of the finitely many coefficients
gives CS29 on all of $W$.  Combining it with the already verified
edge, metric, scalar, and independently closed-pair corrections gives
(CT17).  This is stronger than showing that each numerical
$q_s[\Psi(t)]$ has the same ordinary boundary: it also controls
the single and double exterior-variable coefficients.

The ordinary endpoint evaluations of (CT15) are $\Psi_0,\Psi_1$,
so $\operatorname{ev}_i\widetilde q=q_s[\Psi_i]$.  Each single
exterior coefficient in any derivative polynomial in
$\widetilde\Psi$ contains a factor $\rho'(t)$.  One sees this
directly from the first derivatives and from
\[
 \widetilde D\widetilde\Psi
  =D\Psi_0+\rho DH
   +\epsilon'\rho' d''H-\epsilon''\rho' d'H
   +\epsilon'\epsilon''\rho''H.
 \tag{CT19}
\]
The remaining formulas contain only finite multiplication and full
differentiation; equivalently, differentiating once in the parameter
always leaves a factor $\rho'$, whereas two parameter differentials
belong to the double exterior coefficient.  No double exterior term
can turn into a single exterior term.  Since $\rho'(0)=\rho'(1)=0$,
all single exterior coefficients of $\widetilde q$ vanish at both
endpoints.  In particular (CT6) holds.  Applying (CT8) to (CT17)
proves (CT14).
\end{proof}

\paragraph{Entire-presentation and field checks.}
All operations in (CT13)--(CT19) use the original units and all their
real absolute-value and cap walls, the original complete metric,
and one finite common refinement of the supplied whole presentation.
They commute with restriction and refinement, including presentations
with non-product units such as $(T_1,T_2,T_1+T_2)$ wherever all three
are units.  The argument is pointwise over completed residue fields,
not restricted to rational points.  Passive factors and moving
punctures add no hypothesis: equality is verified on the ordinary
opens (CT18), not on a rational cover substituted for actual opens.
No field extension, higher actual descent, domain enlargement, or
extra growth assumption occurs.

\subsection{What this removes, and what it does not remove, in the cube}

For the raw six-face CS cube QH10, (CT14) implies that replacing the
positive cap constant on any one face changes the cube by an explicitly
specified $\partial_3$ boundary.  A finite list of replacements is
handled by the signed sum of the corresponding $H_{01}$.  Thus the
existence of a relative witness for QH10 is independent of those cap
constants; arbitrary positive cap constants are not an additional
obstruction.

This conclusion does not compare distinct passive blocks.  In QH10,
for example, $q_{a,b}(f)$ and $q_{a,b}(Cf)$ normalize by
$(f_3,\ldots,f_r)$ and its $C$-changed tuple respectively.  Their
uncapped norm potentials need not agree on the passive-active open.
Then $H=0$ on $W_0$ in (CT18) is false, the formal mixed curvature
can survive there, and the passive projective nilpotence used in
the proof of (CT17) is no longer justified for the interpolation.
The finite parameter criterion is therefore not an automatic
solution of the raw cube.

Likewise, six faces would cancel term by term if a single actual
normalization were known to be simultaneously legal for all of them,
but the existence of that common normalization is not supplied here.
Replacing each of six different caps by an arbitrary shared function
without proving the full CS29 identity is not permitted.

No new original-$U$ Green representative is asserted by this module:
it proves a comparison of specified overlap triples.  Applying the
frozen RAG/RSP bridge still requires its whole-overlap relative data
and, for a general finite cover, the remaining independent coherence
hypotheses.  The statement is not a reproof of the frozen CS existence
class and is not a claim that a failure of the cube would disprove G.

\paragraph{Exact dependencies.}
The only frozen mathematical inputs used are the actual differential
algebra and carrier interfaces already used in the exact recovered
R10 CS module, its CS10--CS36 formulas, exact MT8--MT21, the actual
pair homotopy CPT, and the explicit triple/quadruple conventions
QH1--QH3.  Their paths and SHA256 hashes are recorded in the companion
dependency JSON.  The FIV five-potential identity is not invoked to
identify a raw capped-CS triple with a PPC tetrahedron.  No unrecovered
R10 source is used.
